% =============================================================================
% MIN-REPORT-RL: A Minimum-Reportable-Stack Standard for
%                Reinforcement-Learning Post-Training of LLMs
% =============================================================================
% Pillar 4 of the ZVF Program. POSITION PAPER (NeurIPS/ICML Position Track +
% MLRC / ICLR Reproducibility Track).
%
% This is a standalone draft. It is self-contained and should compile with a
% plain `article` class (no external .sty or .bib required: citations are
% rendered as inline TODO keys via a redefined \cite, and the bibliography is
% a manual thebibliography stub). When promoting to a venue template, swap the
% preamble for neurips_2026.sty + \usepackage[numbers]{natbib} and replace the
% TODO \cite keys / thebibliography stub with a real .bib.
% =============================================================================

\documentclass[11pt]{article}

\usepackage[utf8]{inputenc}
\usepackage[T1]{fontenc}
\usepackage[margin=1in]{geometry}
\usepackage{hyperref}
\usepackage{url}
\usepackage{booktabs}
\usepackage{amsmath}
\usepackage{amssymb}
\usepackage{xcolor}
\usepackage{enumitem}
\usepackage{microtype}
\usepackage{graphicx} % for \resizebox on wide tables
\usepackage{listings} % for JSON schema listings
\lstset{basicstyle=\small\ttfamily,frame=single,columns=fullflexible,breaklines=true}

% ---------------------------------------------------------------------------
% TODO-citation shim. We do NOT invent bib entries with fake authors/years.
% Every \cite below uses a clearly-TODO key and renders in red so a human can
% find and replace it with a real reference + .bib entry before submission.
% ---------------------------------------------------------------------------
\renewcommand{\cite}[1]{\textcolor{red}{[\textsc{cite:}\,#1]}}
\newcommand{\todo}[1]{\textcolor{red}{\textbf{[TODO:\ #1]}}}
\newcommand{\zvf}{\textsc{Zvf}}
\newcommand{\gu}{\textsc{Gu}}
\newcommand{\minreport}{\textsc{Min-Report-RL}}
\newcommand{\grpoReg}{\textsc{GRPO-Reg}}

\title{\minreport{}: Reporting the Stack, Not the Label\\
\large A Community Minimum-Reportable-Stack Standard and
Reproducibility-Audit Protocol\\for Group-Relative RL Post-Training of LLMs}

\author{Arvind C R \and the ZVF Program\thanks{Pillar~4 of the ZVF Program.
Author list and affiliations \todo{finalize author block and affiliations}.}}

\date{\today}

\begin{document}
\maketitle

% =============================================================================
\begin{abstract}
% =============================================================================
Reinforcement-learning post-training of large language models is reported as
though an algorithm label---PPO, GRPO, DPO, and the growing GRPO family
(DAPO, GSPO, Dr.GRPO, MAD-GRPO, GRESO, EDGE-GRPO, DARS, TreePo,
\ldots)---fixes the experiment. It does not.
In a controlled audit of a group-relative RL runner, nominally identical
``GRPO'' configurations (same model, group size, learning rate, dataset,
seed, and step budget) produced last-10 training reward of $84.4\%$~\todo{trace to v1 audit citation} on one
backend and $5.0\%$~\todo{trace to v1 audit citation} on another---a $\sim 17\times$~\todo{trace to v1 audit citation} gap with \emph{no visible
hyperparameter difference}. The label was held constant; the \emph{stack} was
not. We argue that this stack-conditioning is not an exotic edge case but the
default state of the field, and that the literature's current reporting norms
make it structurally impossible to tell whether a claimed algorithmic gain
survives a change of backend, sampler, reference-policy/KL handling, LoRA
configuration, or reward parser.

This position paper proposes \minreport{}: an eight-item
\emph{minimum-reportable-stack} checklist that every GRPO-family paper should
satisfy, where each item is included precisely because it is a documented
lever that can flip a head-to-head comparison. We then propose a concrete
reproducibility-audit protocol---re-implement DAPO, GSPO, Dr.GRPO, and
MAD-GRPO inside \emph{one} controlled stack and report which claimed gains
survive---and we specify the audit's experimental design, with results tables
whose cells are explicit placeholders to be filled from the audit corpus. We
close with an adoption path for getting TRL, verl, and OpenRLHF to log the
\minreport{} fields by default, and with responses to the obvious objections.
We do \emph{not} claim that any specific GRPO variant is overstated; we claim
that, under current reporting, the field cannot \emph{know}, and that this is
fixable with a few lines of telemetry.
\end{abstract}

% =============================================================================
\section{Introduction: The Telemetry and Reporting Gap}
\label{sec:intro}
% =============================================================================

\noindent\emph{Scope note (2026-07-11 canonicalization).} This document is
the condensed community-position statement of the \minreport{} standard.
The full-length treatment --- the measured exhibits, live-corpus coupling
analysis, threat model with flip-risk grades, toolchain, and verifier
implementation --- is the Pillar-5 paper \emph{``Report the Stack, Not the
Label: RL-for-LLM Results Are Stack-Conditioned''}
(\texttt{platform\_hybrid/paper/paper\_P5\_minreport}). The two documents
share the same eight-item standard; the Pillar-5 paper is canonical for
evidence, this one for the community-facing statement.

A reader of the 2023--2026 RL-for-LLM literature could be forgiven for
believing that ``we trained with GRPO'' is a complete experimental
description. It is not. The same three-letter label is attached to runs that
differ in their loss form (is there a PPO ratio? a clip? a completion-only
token mask?), their reference-policy and KL handling (frozen reference? KL
penalty in the loss, or in the reward, or absent?), their sampler and backend
(vLLM vs.\ a managed inference API; bf16 vs.\ fp32 logits), their group-size
schedule, and the parser that converts a generation into a scalar reward. Each
of these is, individually, sufficient to move a result. Together they form a
\emph{treatment} that the algorithm label does not name.

This is the LLM-post-training analogue of the deep-RL reproducibility crisis
documented a decade ago \cite{henderson2018deeprl, islam2017reproducibility},
now recurring one abstraction level up. The earlier crisis was about code-level
and hyperparameter variance within a single algorithm; the present one is
worse, because the ``stack'' spans a managed API whose loss we cannot inspect,
a tokenizer and chat template that silently change the prompt, a sampler whose
numerics differ from the trainer's, and a reward parser that can reward a
format artifact instead of a correct answer.

\paragraph{The concrete trigger.}
In an audit of a critic-free group-relative runner \cite{zvfaudit2026}, we
attempted a deliberately boring comparison: hold the \emph{visible} GRPO
configuration fixed---Qwen3-8B, group size $G=8$, learning rate $10^{-5}$,
GSM8K, 30 steps, seed~42---and swap only the backend. A managed runner reached
$84.4\%$~\todo{trace to v1 audit citation} last-10 training reward; TRL on an H100 reached $5.0\%$~\todo{trace to v1 audit citation}. No visible
hyperparameter explains the gap. What differs is everything the label omits:
checkpoint and tokenizer identity, prompt construction, sampler behavior, loss
masking, KL/reference handling, optimizer defaults, LoRA target modules,
precision, rollout plumbing, checkpoint selection, and the evaluator. The
comparison is not evidence that one backend is better; it is evidence that
``the GRPO config'' was \emph{under-specified}. If a backend swap can move a
result by $17\times$~\todo{trace to v1 audit citation} with no visible knob change, then a head-to-head between
two GRPO \emph{variants}---each implemented in its own stack---tells us almost
nothing about the variants.

\paragraph{The family is growing faster than its reporting.}
The TMLR survey of agentic RL \cite{agenticrlsurvey2026} now catalogs, in a
single comparison table, more than twenty GRPO-family variants---and a
striking fraction of them (DAPO's dynamic sampling, GRESO's pre-rollout
filtering, EDGE-GRPO's \emph{advantage collapse} mitigation, DARS's
difficulty-aware reallocation) intervene on the same underlying lever: the
fraction of groups whose identical rewards contribute zero gradient, i.e.\
the quantity \zvf{} measures. Not one of the cataloged variant papers
reports a per-step dead-group trajectory or the stack fields below, so
their relative standings are, under current norms, unattributable. The
same survey's mechanistic synthesis supplies a second indictment of current
reporting: roughly two-thirds of the RL-for-reasoning studies it reviews
report only pass@1, although pass@$k$ curves are what distinguish a model
whose solution \emph{distribution was sharpened} from one whose solution
\emph{support expanded}. A variant can ``win'' at pass@1 by concentrating
probability on already-reachable solutions while strictly shrinking the
set of problems it can solve at pass@32. Under pass@1-only reporting, these
two outcomes---one an optimization of sampling, the other a capability
change---are indistinguishable. This motivates the eighth item of the
standard.

\paragraph{Why this is a reporting problem, not just a science problem.}
The gap above is not a bug to be fixed by a better trainer; it is information
that is routinely \emph{not logged}. Most papers do not report whether the KL
term was in the loss or the reward, whether the token mask was completion-only,
what the sampler precision was, or what the per-step group-size schedule looked
like. They also rarely report the per-step optimization telemetry---in our
framework, the Zero-Variance Fraction (\zvf{}) and Gradient Utilization
(\gu{})---that would reveal whether a run had any usable learning signal at all
\cite{zvfaudit2026}. A field cannot audit what it does not record. The fix is
cheap: a small, fixed set of fields, logged by default.

\paragraph{Scope and provenance of the worked numbers.}
The concrete numbers in this paper used as motivating examples
($84.4\%$/$\,5.0\%$/$\,17\times$ in the audit;
$61.6$--$89.6\%$ prompt-token loss magnitude;
$82.0\%\to83.3\%$ held-out control at $p=0.26$;
$95.0$--$98.1\%$/$\,95.0\%$ Llama-3.3-70B run)
are \emph{inherited from the v1 audit paper that motivated this
position piece} \todo{cite v1 audit paper, once assigned} and are
marked inline with \todo{trace to v1 audit citation} wherever they
appear. A reader who wants to verify the numbers should consult that
paper; this manuscript presents them as illustrative of the
\emph{kind} of stack-driven flip, not as re-derivable from a fresh
run. Once the v1 audit paper is in citation scope, the inline
\todo{trace} markers are replaced with real \texttt{\textbackslash cite\{zvfaudit\}}
keys and the corresponding bibliography entry below is filled in.
We do \emph{not} claim that any specific number here is independently
verifiable against the present manuscript's experimental record.

\paragraph{Contributions.}
This paper makes three:
\begin{enumerate}[leftmargin=1.6em, itemsep=2pt]
  \item \textbf{The stack-conditioning thesis, sharpened into a reporting
        standard.} We restate and extend the audit finding that algorithm
        labels are under-specified treatments (\S\ref{sec:stack}), and convert
        it into \minreport{}: an eight-item minimum-reportable-stack checklist
        in which every item is justified by a mechanism through which it can
        flip a comparison (\S\ref{sec:standard}).
  \item \textbf{A controlled reproducibility-audit protocol.} We specify a
        single-stack re-implementation of DAPO, GSPO, Dr.GRPO, and MAD-GRPO,
        with a pre-registered analysis that reports which claimed gains survive
        when the stack is held fixed (\S\ref{sec:audit}). Result cells are
        explicit placeholders to be filled from the audit corpus.
  \item \textbf{An adoption path.} We describe concrete, low-friction changes
        to TRL, verl, and OpenRLHF so that the \minreport{} fields are emitted
        by default, plus a venue checklist (\S\ref{sec:adoption}).
\end{enumerate}

% =============================================================================
\section{The Stack-Conditioning Problem}
\label{sec:stack}
% =============================================================================

\paragraph{Thesis.}
\emph{Algorithm rankings in RL post-training are stack-conditioned.} A claim of
the form ``method $A$ beats method $B$'' is, in current practice, a claim about
two \emph{full stacks} $S_A$ and $S_B$ that happen to be tagged $A$ and $B$. The
quantities that determine the outcome---backend and sampler numerics, the exact
loss form (ratio, clip, token mask), reference-policy and KL handling, LoRA
target modules and rank, precision, the reward parser, the group-size schedule,
checkpoint selection, and the evaluation harness---are co-determinants of the
result, and they are typically \emph{confounded with the label}. In a one-way
variance decomposition over a heterogeneous run corpus, the framework/library
factor and the training-algorithm factor each explain roughly half of the
variance in the outcome ($\eta^2 \approx 0.55$ for both), but these factors are
themselves confounded with model, backend, reward, hardware, and task
\cite{zvfaudit2026}. ``Algorithm'' and ``stack'' are not separable in the
existing literature because nobody reports the stack.

\paragraph{A comparison that flips.}
Consider the minimal worked example from the audit. The \emph{visible} GRPO
configuration is matched exactly:
\begin{center}
\small
\begin{tabular}{@{}lll@{}}
\toprule
\textbf{Visible config (held constant)} & \textbf{Backend $S_1$} & \textbf{Backend $S_2$}\\
\midrule
Model & Qwen3-8B & Qwen3-8B \\
Group size $G$ & 8 & 8 \\
Learning rate & $10^{-5}$ & $10^{-5}$ \\
Dataset / split & GSM8K-500 & GSM8K-500 \\
Steps & 30 & 30 \\
Seed & 42 & 42 \\
\midrule
\textbf{Last-10 training reward} & $\mathbf{84.4\%}$~\todo{trace} & $\mathbf{5.0\%}$~\todo{trace} \\
\bottomrule
\end{tabular}
\end{center}
A na\"ive reader concludes ``$S_1$'s GRPO is $17\times$~\todo{trace to v1 audit citation} better.'' The correct
reading is that the two rows are \emph{different treatments wearing the same
label}: the loss form, token mask, KL/reference handling, sampler precision,
LoRA targets, optimizer defaults, and reward parser were never held fixed
because they were never reported, hence never matched. Flip any one of them and
the ranking can invert. This is the entire problem in one table: the label is
constant, the conclusion is determined by the unreported stack.

\paragraph{Why the GRPO family makes this acute.}
The methods this paper targets---DAPO, GSPO, Dr.GRPO, MAD-GRPO, and the broader
variance-mitigation line (AERO, CPPO, NGRPO, Scaf-GRPO)
\cite{dapo2025, gspo2025, drgrpo2025, madgrpo2025, aero2024, cppo2024,
ngrpo2025, scafgrpo2025}---are typically defined as \emph{small deltas} on a
base GRPO loop: a changed advantage normalization, a clip-bound tweak, a token
mask, a length penalty, an exploration bonus, an adaptive group size. The size
of the claimed improvement is frequently \emph{smaller} than the stack effect
demonstrated above. When the treatment effect you are trying to measure is
$2$--$5$ points and the nuisance effect of an unreported stack difference is
tens of points, an uncontrolled head-to-head is not measuring the method. The
GRPO family is exactly the regime in which stack-conditioning is most likely to
masquerade as algorithmic progress.

% =============================================================================
\section{The \minreport{} Standard}
\label{sec:standard}
% =============================================================================

\minreport{} is a minimum-reportable-stack: the smallest set of fields such
that, if two GRPO-family papers both report them, a reader can tell whether
their comparison is confounded. Each item below is included \emph{because there
is a known mechanism by which it can flip a comparison}; an item that could not
change a ranking would not earn its place on a minimum list.

\paragraph{1. Loss form.}
\emph{Report:} whether the update uses a PPO-style importance ratio
$w_{i,t}=\pi_\theta/\pi_{\theta_{\text{old}}}$; whether and how it is clipped
(and the clip bounds, including asymmetric DAPO-style ``clip-higher'');
whether the token mask is completion-only or whole-sequence; and whether
advantages are normalized per-group, per-batch, or with a running estimate.
\emph{Why it can flip:} the choice of token mask alone changes which tokens
carry gradient. In one diagnostic, $61.6$--$89.6\%$~\todo{trace to v1 audit citation} of the full-sequence loss
magnitude came from \emph{prompt} tokens rather than completion tokens
\cite{zvfaudit2026}; a whole-sequence mask and a completion-only mask are
therefore different objectives sharing a name. Dr.GRPO's contribution is
precisely a change to length/normalization in the loss \cite{drgrpo2025}; GSPO
changes the importance-sampling granularity \cite{gspo2025}. If the baseline's
loss form is unreported, the variant's gain is unattributable. The stake is not
hypothetical: in our own companion panel, a documented-but-unwired loss flag
left six ``GRPO vs.\ Dr.GRPO'' arms silently training the identical objective,
and no reward, length, or ZVF trace revealed it --- only reading the runner
did. Loss form must be \emph{reported and verified}, not inferred from
plausible-looking outputs.

\paragraph{2. Reference policy and KL handling.}
\emph{Report:} whether a frozen reference policy is retained; whether the KL
term sits in the \emph{loss} or is folded into the \emph{reward}; the KL
coefficient and its schedule; and whether KL is estimated forward or reverse,
per-token or per-sequence. \emph{Why it can flip:} KL placement changes the
effective objective and the steady-state exploration. A runner that drops the
reference policy entirely (one optimizer step per rollout, no frozen reference)
is not doing the same thing as a KL-regularized GRPO, even at matched learning
rate \cite{zvfaudit2026}. Methods that add a second anchor (e.g.\ a dual-KL/SFT
anchor) change steady-state \zvf{} and must be recalibrated, not compared
blind \cite{dar2024}.

\paragraph{3. Sampler, backend, and precision.}
\emph{Report:} the rollout engine (vLLM, SGLang, a managed API, the trainer's
own \texttt{generate}; decoding parameters (temperature, top-$p$,
\texttt{max\_tokens}; logit precision in the sampler vs.\ the trainer (bf16 /
fp16 / fp32); and whether sampler and trainer share the same tokenizer and chat
template. \emph{Why it can flip:} the sampler defines the rollout distribution
that the group-relative update consumes. A managed, closed-source runner and an
open vLLM path can yield order-of-magnitude different training-reward
trajectories under identical visible configs---this is the mechanism behind the
$17\times$ gap of \S\ref{sec:stack}. Sampler precision shifts the probability
of mixed-reward groups, and hence the available gradient.

\paragraph{4. Per-step \zvf{} and \gu{} trajectory.}
\emph{Report:} the per-step Zero-Variance Fraction \zvf{} (fraction of prompts
whose $G$ completions all receive identical reward, contributing zero gradient)
and Gradient Utilization $\gu{}=1-\zvf{}$, logged for every training step.
\emph{Why it can flip:} \zvf{}/\gu{} is the telemetry that reveals \emph{whether
a run had any learning signal at all}. The mixed-group probability
$P(\text{usable}) \approx \frac{1}{N}\sum_x\!\left[1-(1-p_x)^G-p_x^G\right]$
shows that a group-relative update only has signal when sampled groups contain
reward diversity \cite{zvfaudit2026}. A method can ``win'' a comparison purely
because its stack happened to produce lower \zvf{} (more usable groups), not
because its algorithmic idea is better. Without the \zvf{}/\gu{} trajectory, a
collapsed run and a saturated run are indistinguishable from their final reward.
A fixed first-five-step rule (\zvf{}$\,\ge 80\%$ with reward $\le 5\%$) is a
cheap collapse triage; reporting the trajectory lets a reader apply it.

\paragraph{5. Group-size schedule (fixed or adaptive).}
\emph{Report:} the group size $G$ at every step, and the rule that changes it
if it is adaptive (e.g.\ AERO-style doubling/halving on a rolling \zvf{}
estimate). \emph{Why it can flip:} $G$ directly controls the mixed-group
probability above and is one of the strongest single knobs we measured---an
ablation on Qwen3-8B/GSM8K moved last-10 reward across a wide band as $G$ went
$2\!\to\!4\!\to\!8\!\to\!16$ \cite{zvfaudit2026}. An adaptive-$G$ method (AERO)
compared against a fixed-$G$ baseline is partly being credited for spending
more compute on hard prompts, an effect that must be separated from the
algorithmic claim \cite{aero2024}.

\paragraph{6. Held-out split distinct from the reward environment.}
\emph{Report:} a held-out evaluation slice that is \emph{disjoint} from the
training prompts and scored by a harness, with sample size and a confidence
interval, reported \emph{separately} from online training reward. \emph{Why it
can flip:} online training reward is dynamics evidence, not capability
evidence. In the audit, a clean paired control improved only $82.0\%\to83.3\%$~\todo{trace to v1 audit citation}
($p=0.26$)~\todo{trace to v1 audit citation} on held-out GSM8K despite training reward near saturation; and
selecting checkpoints by training reward produced an $87$--$95\%$ ``capability''
band that is a selection artifact, not a causal lift \cite{zvfaudit2026}. A
variant that reports only training reward can show a large apparent gain that
vanishes held-out. Four Llama-3.3-70B seeds ranged $95.0$--$98.1\%$~\todo{trace to v1 audit citation} on training
last-10 yet all landed on $95.0\%$~\todo{trace to v1 audit citation} held-out---last-10 variance was
sampling noise, not generalization.

\paragraph{7. Decontamination probe results.}
\emph{Report:} the outcome of an explicit train/test contamination check
(n-gram or embedding overlap between training prompts and the held-out /
benchmark slice), and the parser's behavior on adversarial format-only inputs.
\emph{Why it can flip:} verifiable rewards do not eliminate reward hacking;
parser artifacts, format shortcuts, length effects, and train-prompt overfitting
can all inflate a reward without capability gain \cite{zvfaudit2026}. If the
benchmark slice overlaps the reward environment, a ``gain'' may be memorization;
if the parser rewards a format token, a ``gain'' may be a shortcut. A
contamination/parser probe is the minimum evidence that the reported number
measures the thing it claims to measure.

\paragraph{8. Pass@$k$ curves alongside pass@1.}
\emph{Report:} held-out pass@$k$ at $k \in \{1, 8, 32\}$ (or a stated
subset with justification), with the sampling temperature and completion
budget used for the estimate. \emph{Why it can flip:} pass@1 conflates two
different improvements---sharpening the output distribution over solutions
the base model could already reach, and expanding the set of reachable
solutions. The agentic-RL survey's synthesis \cite{agenticrlsurvey2026}
finds that roughly two-thirds of RL-for-reasoning papers report only
pass@1, while the studies that do report pass@$k$ frontiers repeatedly find
base models matching or overtaking their RL-tuned counterparts at large
$k$---i.e., the ``gain'' was distribution sharpening. A GRPO-variant
comparison scored at pass@1 alone can therefore rank the method that
sharpens above the method that generalizes, and the ranking can invert at
$k=32$. Because group-relative training acts directly on the sampling
distribution, this is not an exotic failure mode for the GRPO family; it is
the expected one.

\paragraph{Summary table.}
Table~\ref{tab:checklist} collects the standard. The companion
\texttt{CHECKLIST.md} provides a fillable appendix template.

\begin{table}[t]
\centering
\caption{The \minreport{} minimum-reportable-stack. Each item is on the list
because it is a documented lever that can flip a head-to-head comparison.}
\label{tab:checklist}
\small
\begin{tabular}{@{}rp{0.30\linewidth}p{0.50\linewidth}@{}}
\toprule
\# & \textbf{Field} & \textbf{Flip mechanism (why it is mandatory)}\\
\midrule
1 & Loss form (ratio / clip / token mask / advantage norm.) & Token mask reassigns gradient across prompt vs.\ completion tokens; clip and norm.\ are the actual variant deltas.\\
2 & Reference policy + KL handling & KL placement (loss vs.\ reward) and a frozen-vs-absent reference change the objective and exploration.\\
3 & Sampler / backend / precision & The sampler defines the rollout distribution; backend/precision drove the $17\times$~\todo{trace to v1 audit citation} matched-config gap.\\
4 & Per-step \zvf{} / \gu{} trajectory & Reveals whether the run had usable learning signal; separates collapse from saturation.\\
5 & Group-size schedule (fixed / adaptive) & $G$ sets the mixed-group probability; adaptive-$G$ confounds compute with algorithm.\\
6 & Held-out split $\neq$ reward environment & Training reward is not capability; selection-by-reward inflates apparent gains.\\
7 & Decontamination probe results & Verifiable rewards still admit parser/format/length/overlap hacking.\\
8 & Pass@$k$ curves alongside pass@1 & Pass@1-only scoring cannot distinguish distribution sharpening from capability gain; rankings can invert at large $k$.\\
\bottomrule
\end{tabular}
\end{table}

\subsection{Measured stack effects (new evidence, 2026-07)}
\label{sec:measured}

Three of the standard's items now carry \emph{quantified} first-party
evidence from the companion experiments suite
(\texttt{zvf-program/experiments-next/results/}), replacing argument by
anecdote.

\paragraph{Item 3, measured: the sampler moves frontier metrics.}
Holding model, prompts (200 GSM8K test problems), temperature, seed, and
parser byte-identical, and swapping \emph{only} the sampling backend:
\begin{center}
\begin{tabular}{@{}lccc@{}}
\toprule
Backend & pass@1 & pass@8 & pass@32 \\
\midrule
Managed (Tinker) sampler & 30.4\% & 79.7\% & 91.0\% \\
vLLM 0.24 on A100 & 30.6\% & 77.8\% & 89.0\% \\
\bottomrule
\end{tabular}
\end{center}
Point estimates are backend-robust ($+0.2$ at pass@1) while the pass@32
frontier shifts by $2.0$ points and individual problems move by up to
$|11|/32$ solutions. Any cross-paper comparison of frontier metrics that
does not report the sampler is therefore confounded at the
$\sim$2-point scale --- the same scale as many claimed algorithmic gains.

\paragraph{Item 1, measured: silent export/loader incompatibilities.}
Evaluating LoRA adapters trained on a managed stack inside an open
inference stack required reverse-engineering two undocumented format
deltas: \texttt{target\_modules="all-linear"} (PEFT shorthand the open
loader rejects) and the unembedding trained under
\texttt{model.unembed\_tokens} (the open loader accepts only
\texttt{lm\_head}). Without remapping, \emph{every} cross-stack evaluation
of these 21 public adapters fails --- loudly in our harness, but a wrapper
that silently skipped the adapter would have produced base-model numbers
labeled as post-RL results. This is precisely the class of confound item 1
exists to surface.

\paragraph{Item 8, exemplified: the panel pass@1 alone would misreport.}
A matched base-vs-post-RL panel (four adapters trained at
$G\in\{2,4,16,32\}$; one stack; same problems) shows pass@1 deltas
scattered from $-0.5$ to $+1.9$ points --- under pass@1-only reporting,
the $G{=}2$ arm would be called a \emph{regression}. The pass@32 frontier
tells the opposite, consistent story: all four arms improve by $+1.5$ to
$+2.5$ points. Neither reading reaches significance at a single training
seed (the deltas sit within the measured sampler and eval-seed noise
bands above), which is itself the point: \textbf{without pass@$k$ curves
and stated noise bands, this panel supports any narrative a reader
prefers.} With them, it supports exactly one claim: a small,
directionally consistent frontier shift requiring multi-seed replication.

% =============================================================================
\section{Proposed Reproducibility Audit}
\label{sec:audit}
% =============================================================================

The \minreport{} standard tells authors what to report. The audit tells the
community what to \emph{check}: \emph{which claimed GRPO-family gains survive
when the stack is held fixed?}

\paragraph{Design.}
Re-implement four prominent GRPO variants---DAPO, GSPO, Dr.GRPO, and
MAD-GRPO \cite{dapo2025, gspo2025, drgrpo2025, madgrpo2025}---as
\emph{minimal configuration overrides on a single shared trainer}. Everything
in the \minreport{} list except the variant's defining delta is held identical
across arms: the same base checkpoint and tokenizer, the same vLLM sampler and
precision, the same reference-policy/KL handling, the same LoRA targets and
rank, the same optimizer, the same reward parser, the same group-size schedule
(where the variant does not itself change it), the same held-out slice, and the
same evaluation harness and seed sweep. Only the named hook differs (loss form,
advantage normalization, clip rule, masking, or group-sizing). This is the same
methodology already validated for the variance-mitigation head-to-head, where
AERO/CPPO/NGRPO/Scaf-GRPO were each implemented as one hook on a shared GRPO
trainer with tokenizer, sampler, optimizer, LoRA adapters, evaluation harness,
and seed sweep held identical \cite{zvfaudit2026}.

\paragraph{Two readings per variant.}
For each variant we report (i) its result \emph{as published} (its own stack,
its own numbers) and (ii) its result \emph{in the controlled stack}. The
quantity of interest is the \emph{survival} of the claimed gain: the
controlled-stack delta relative to the shared GRPO baseline, with a confidence
interval and a survival verdict (survives / shrinks / disappears / reverses).

\paragraph{Pre-registration.}
The arm list, the shared stack specification, the held-out slice, the seed set,
and the survival thresholds are fixed before any controlled run is scored, to
avoid the selection-by-reward over-optimism the audit itself warns against
\cite{zvfaudit2026}. Every arm reports the full \minreport{} block.

\paragraph{Primary results table.}
Table~\ref{tab:audit} is the headline deliverable; all numeric cells are
placeholders to be filled from the audit corpus.

\begin{table}[t]
\centering
\caption{\textbf{Controlled single-stack audit of GRPO-family variants.}
Each variant is a minimal hook on one shared trainer (\S\ref{sec:audit}); the
full \minreport{} block is identical across arms except the variant's defining
delta. ``Published $\Delta$'' is the gain the original paper reports on its own
stack; ``Controlled $\Delta$'' is the gain in the shared stack vs.\ the shared
GRPO baseline; ``Survives?'' is the pre-registered verdict. All cells are
\todo{fill from audit corpus}.}
\label{tab:audit}
\small
\resizebox{\textwidth}{!}{%
\begin{tabular}{@{}lcccccc@{}}
\toprule
\textbf{Variant} & \textbf{Defining delta} & \textbf{Published $\Delta$} &
\textbf{Controlled last-10} & \textbf{Controlled held-out} &
\textbf{Controlled $\Delta$ (95\% CI)} & \textbf{Survives?}\\
\midrule
GRPO (baseline) & --- & --- & \todo{} & \todo{} & $0$ (ref.) & --- \\
DAPO     & clip-higher + dyn.\ sampling & \todo{} & \todo{} & \todo{} & \todo{} & \todo{} \\
GSPO     & sequence-level IS ratio      & \todo{} & \todo{} & \todo{} & \todo{} & \todo{} \\
Dr.GRPO  & length/normalization fix     & \todo{} & \todo{} & \todo{} & \todo{} & \todo{} \\
MAD-GRPO & multi-agent / diversity term & \todo{} & \todo{} & \todo{} & \todo{} & \todo{} \\
\bottomrule
\end{tabular}}
\end{table}

\paragraph{Telemetry table.}
Table~\ref{tab:audit_telemetry} reports the \minreport{} item-4 telemetry for
each arm, which is the mechanism by which a ``gain'' might be a stack effect
rather than an algorithmic one (e.g.\ a variant that simply lowers \zvf{}).

\begin{table}[t]
\centering
\caption{\textbf{Per-arm \minreport{} telemetry.} \zvf{}/\gu{} and collapse
behavior under the shared stack, to attribute any controlled gain to algorithm
vs.\ usable-signal availability. All cells are \todo{fill from audit corpus}.}
\label{tab:audit_telemetry}
\small
\resizebox{\textwidth}{!}{%
\begin{tabular}{@{}lccccc@{}}
\toprule
\textbf{Variant} & \textbf{Mean \zvf{} @ step 25} & \textbf{Mean \gu{} @ step 25} &
\textbf{Collapse rate (seeds)} & \textbf{Time-to-collapse (median steps)} &
\textbf{Group-size schedule}\\
\midrule
GRPO (baseline) & \todo{} & \todo{} & \todo{} & \todo{} & fixed $G$ \\
DAPO     & \todo{} & \todo{} & \todo{} & \todo{} & \todo{} \\
GSPO     & \todo{} & \todo{} & \todo{} & \todo{} & \todo{} \\
Dr.GRPO  & \todo{} & \todo{} & \todo{} & \todo{} & \todo{} \\
MAD-GRPO & \todo{} & \todo{} & \todo{} & \todo{} & \todo{} \\
\bottomrule
\end{tabular}}
\end{table}

\paragraph{Statistical protocol.}
Because the audit's independent unit is a seed (not a per-step reward, which is
autocorrelated within a trace), inferential claims are made only across seeds.
We report mean $\pm$ standard error and bootstrap 95\% CIs across the seed
sweep, control the false-discovery rate with Benjamini--Hochberg over the
pre-registered set of survival tests, and treat single-seed arms as descriptive
only. This mirrors the statistical discipline of the source audit, which
explicitly refuses to upgrade trace-level descriptive effect sizes into
inferential evidence about distinct seeds, runs, or stacks \cite{zvfaudit2026}.
\todo{Fix the seed count $S$ and the survival threshold once compute is
budgeted; record minimum detectable effect at the chosen $S$.}

\paragraph{Honest scoping.}
The audit's deliverable is \emph{not} ``method $X$ is fake.'' It is a survival
map: for each variant, the fraction of its claimed gain that persists when the
stack is the same. Some gains will survive (they are real algorithmic
contributions); some will shrink (they were partly stack); some may disappear
(they were stack). All three outcomes are informative, and reporting them is
the contribution.

% =============================================================================
\section{Adoption Path}
\label{sec:adoption}
% =============================================================================

A standard nobody can comply with cheaply will not be adopted. \minreport{} is
designed so that the marginal cost to an author is close to zero \emph{if the
trainers log it by default}.

\paragraph{Make the trainers emit it.}
The eight fields are already computable inside every major open trainer. The
marginal work is to write them down once, in a stable schema, and to compute
\zvf{}/\gu{} from the grouped reward tensor that GRPO already materializes.

\subparagraph{A concrete TRL plugin: \texttt{trl-min-report-rl}.}
We sketch a small TRL extension that makes adoption automatic. For every
\texttt{GRPOTrainer.train()} call it emits a run-start JSON block containing the
eight \minreport{} fields, then logs per-step \zvf{}/\gu{} from the actual
grouped rewards used by GRPO. The package shape is:
\begin{verbatim}
trl_min_report_rl/
  __init__.py
  config.py              # MinReportRLConfig dataclass + CLI flags
  callback.py            # run-start manifest + finalization
  grpo_mixin.py          # reward-tensor hook
  emitter.py             # JSONL + stdout + logger adapters
  schema/
    grpo_run_start.v0.1.json
    grpo_step.v0.1.json
\end{verbatim}
The public API is a single extra argument to \texttt{GRPOConfig}:
\begin{verbatim}
args = GRPOConfig(..., report_min_report_rl=True)
trainer = MinReportGRPOTrainer(
    model=model, args=args, reward_funcs=[reward_fn],
    train_dataset=train_ds, eval_dataset=heldout_ds,
    processing_class=tokenizer, min_report_rl=MinReportRLConfig(...))
\end{verbatim}
A no-subclass adapter, \texttt{enable\_min\_report\_rl(trainer, config)}, is
provided for compatibility with existing code.

The critical hook is \texttt{on\_grpo\_rewards\_computed}, called after reward
functions return and after distributed gather, but before advantage
normalization. It receives \texttt{rewards\_per\_func}, \texttt{reward\_weights},
and \texttt{num\_generations}, and computes
\[
\text{ZVF}_t = \frac{1}{N}\sum_{x=1}^{N}
  \mathbf{1}\big[\operatorname{Var}(r_{x,1},\dots,r_{x,G})=0\big],
\qquad
\text{GU}_t = 1 - \text{ZVF}_t,
\]
where unusable (all-NaN) groups are counted as zero-variance. The hook logs both
human-readable scalars (\texttt{min\_report\_rl/zvf},
\texttt{min\_report\_rl/gu}, \texttt{min\_report\_rl/usable\_groups}) and a
per-step JSON event.

\subparagraph{Upstream surface.}
The cleanest upstream TRL change adds three elements to \texttt{GRPOConfig}:
\texttt{report\_min\_report\_rl}, \texttt{min\_report\_rl\_strict}, and
\texttt{min\_report\_rl\_output\_dir}; plus two trainer methods
(\texttt{\_emit\_min\_report\_rl\_run\_start} and
\texttt{\_emit\_min\_report\_rl\_step}) and the reward-path hook call. That
is the ``few lines of telemetry'' version of the standard.

\subparagraph{Run-start manifest.}
The first line of \texttt{min\_report\_rl.jsonl} is a JSON object whose
\texttt{min\_reportable\_stack} field contains the eight items. A fragment is
shown in Listing~\ref{lst:runstart}; the full schema is in the
\grpoReg{} catalog paper \cite{grpoRegistry2026}.

\begin{lstlisting}[caption={Fragment of the
  \texttt{min\_report\_rl} run-start manifest.}, label={lst:runstart}]
{
  "event": "min_report_rl.run_start",
  "min_reportable_stack": {
    "1_loss_form": {"status":"known", "loss_type":"dapo", ...},
    "2_reference_policy_kl": {"status":"known", "kl": {"enabled":false}},
    "3_sampler_backend_precision": {"rollout_engine":"vllm", ...},
    "4_zvf_gu_trajectory": {"enabled":true, "basis":"weighted_reward", ...},
    "5_group_size_schedule": {"train_num_generations":8, ...},
    "6_heldout_split": {"heldout_dataset_id":"gsm8k:test:500", ...},
    "7_decontamination_probes": {"ngram_overlap":{"passed":true}, ...}
  }
}
\end{lstlisting}

\subparagraph{Other trainers.}
verl's actor/rollout separation makes sampler, precision, and group-size
schedule first-class; OpenRLHF's reference-model and KL options map directly
onto item~2. Both can emit the same manifest block and per-step telemetry using
the same shared emitter. The schema is trainer-agnostic; only the source of each
field changes.

\paragraph{Enforcement and audit tooling.}
The same manifests feed three tools. \texttt{grpo-stackdiff}
(\S\ref{sec:tooling}) compares two manifests and flags whether their nuisance
differences are large enough to flip a comparison. The \texttt{MIN-REPORT-RL
Auditor} scores an existing paper or repo against the eight-item checklist,
assigns a 0--100 reproducibility-auditable badge, and produces a reviewer-facing
summary \cite{zvfaudit2026}. The \texttt{MRRL-AuditRunner} turns the controlled
single-stack audit of \S\ref{sec:audit} into a pre-registered executable
pipeline: it locks the shared stack, expands each variant as a declared hook,
and computes survival verdicts from paired seed-level held-out deltas
\cite{zvfaudit2026}. Together these tools turn the abstract standard into a
mechanical check that venues and reviewers can run in CI.

\paragraph{Make the venues ask for it.}
Add a \minreport{} block to the reproducibility checklist for the RL-for-LLM
track, analogous to the existing NeurIPS/ICML reproducibility checklists. The
companion \texttt{CHECKLIST.md} is written to be dropped verbatim into a paper's
appendix.

\paragraph{Make it auditable post hoc.}
For already-published work, the same JSON block can be reconstructed from
release artifacts where they exist. The audit of \S\ref{sec:audit} doubles as a
demonstration: it re-derives the \minreport{} block for four variants and shows
the format is sufficient to detect confounded comparisons.

% =============================================================================
\section{Objections and Responses}
\label{sec:objections}
% =============================================================================

\paragraph{``This is just `report your hyperparameters.'\,''}
No. Hyperparameter tables already exist and the $17\times$~\todo{trace to v1 audit citation} gap occurred
\emph{despite} matched visible hyperparameters. \minreport{} names the fields
that hyperparameter tables systematically omit---loss form, KL placement,
sampler precision, per-step \zvf{}/\gu{}, group-size schedule, held-out
disjointness, contamination probe---precisely because each is a demonstrated
flip lever, not a tuning knob.

\paragraph{``Eight items is too many / too few.''}
Each item earns its place by a flip mechanism (\S\ref{sec:standard}); removing
any one reopens a known confound. The list is a \emph{minimum}, not a maximum:
authors may report more. We deliberately stop at eight because that is the set
that, in our audit and in the survey-documented pass@$k$ reporting gap
\cite{agenticrlsurvey2026}, was sufficient to explain the comparisons that
flipped.
\todo{If the audit surfaces a ninth recurring confound, extend the list and
justify it the same way.}

\paragraph{``Re-implementations are unfaithful, so the audit is unfair.''}
The audit reports \emph{both} readings (published and controlled) and frames the
result as survival, not refutation. An unfaithful re-implementation is itself a
finding: if a variant's gain depends on an unstated stack detail that a careful
re-implementer cannot reproduce from the paper, that is a reporting failure the
standard is designed to surface. We will publish every arm's full \minreport{}
block so disagreements are about specific fields, not vibes.

\paragraph{``Closed/managed backends can't comply.''}
Partly true, and that is the point: a closed runner cannot expose its loss form
or KL handling, so results from it should be scoped as ``this platform's
implementation,'' not ``GRPO'' \cite{zvfaudit2026}. The standard makes that
scoping explicit rather than letting a closed-stack number stand in for an
algorithm.

\paragraph{``\zvf{}/\gu{} is your own metric; why mandate it?''}
\zvf{}/\gu{} is one cheap instantiation of a necessary quantity---\emph{did the
run have usable group-relative signal?}---and we accept substitutes. Under
dense process rewards or scaffolded exploration, \zvf{} degenerates and should
be replaced by per-step reward variance or a gradient-norm-variance surrogate
\cite{zvfaudit2026}. Item~4 mandates \emph{a usable-signal trajectory}, of which
\zvf{}/\gu{} is the default; the point is that \emph{some} such telemetry must be
reported.

% =============================================================================
\section{Companion Tooling: \texttt{grpo-stackdiff}}
\label{sec:tooling}
% =============================================================================

A standard is easier to adopt if reviewers can check it mechanically. We sketch
a companion CLI, \texttt{grpo-stackdiff}, that consumes two \minreport{}
run manifests and reports whether their differences are large enough to flip a
comparison. The tool is not a replacement for the audit of
\S\ref{sec:audit}; it is an enforcement layer that tells a reviewer, in
seconds, whether a head-to-head is stack-comparable.

\paragraph{Manifest contract.} Each run is a YAML/JSON block containing the
eight \minreport{} fields, plus run identity (trainer, version, seed, hardware),
model and tokenizer hashes, LoRA configuration, optimizer settings, and results.
The schema is the JSON block proposed in \S\ref{sec:adoption}.

\paragraph{Diff taxonomy.} Every field difference receives three labels:
\begin{itemize}[leftmargin=1.4em, itemsep=2pt]
  \item \textbf{Lever.} One of eight lever groups: L0~run identity/provenance,
    L1~loss form, L2~reference policy and KL handling, L3~sampler/backend/
    precision, L4~usable-signal telemetry (\zvf/\gu{}), L5~group-size
    schedule, L6~held-out evaluation and checkpoint selection, L7~decontamination
    and parser robustness.
  \item \textbf{Diff kind.} EQUAL, COSMETIC, PROVENANCE, PARAMETRIC, SCHEDULE,
    SEMANTIC\_OBJECTIVE, DISTRIBUTIONAL, TELEMETRY, MISSING, OPAQUE, DERIVED, or
    INVALID\_TARGET.
  \item \textbf{Role.} TREATMENT\_DELTA if the user declares it as the
    algorithmic delta being tested; NUISANCE\_DELTA if it is an uncontrolled
    stack difference; COVERAGE\_GAP if the field is missing; INVALIDATOR if the
    comparison target is not common.
\end{itemize}

\paragraph{Flip-risk classes.} The tool compares each nuisance difference to
the reported comparison margin and classifies the overall risk as one of:
\begin{center}
\small
\begin{tabular}{@{}ll@{}}
\toprule
\textbf{Class} & \textbf{Meaning}\\
\midrule
R0 \texttt{same} & No meaningful difference.\\
R1 \texttt{cosmetic} & Version string or comment only.\\
R2 \texttt{small} & Estimated effect $<25\%$ of the comparison margin.\\
R3 \texttt{material} & Effect $25$--$100\%$ of the margin; can shrink a claim.\\
R4 \texttt{flip\_capable} & Effect bound $\	ext{\textgreater=}$ margin; can flip the ranking.\\
R5 \texttt{invalidating} & The comparison target is not common.\\
RU \texttt{unknown} & Missing or opaque evidence; comparison unverifiable.\\
\bottomrule
\end{tabular}
\end{center}

\paragraph{Verdict and CI integration.} The tool emits a deterministic verdict
(\texttt{STACK\_MATCHED}, \texttt{STACK\_MATERIAL}, \texttt{STACK\_CONFOUNDED},
\texttt{UNVERIFIABLE}, or \texttt{INVALID\_COMPARISON}) and an exit code, so a
paper's reproducibility workflow can fail a comparison that is
flip-capable. Example:
\begin{verbatim}
grpo-stackdiff compare dapo.yaml grpo.yaml \
  --metric heldout.exact_match --claim-delta +0.024 \
  --fail-on flip-capable
\end{verbatim}

\paragraph{Role in the \minreport{} ecosystem.} Authors emit the JSON block;
venues ask for it; reviewers run \texttt{grpo-stackdiff}. The cost to an author
is still a few lines of telemetry; the cost to a reviewer is one command. The
tool turns the abstract stack-conditioning thesis of \S\ref{sec:stack} into a
concrete, reproducible check.

% =============================================================================
\section{Conclusion}
\label{sec:conclusion}
% =============================================================================

The RL-for-LLM literature is in a position the deep-RL community has seen
before: a three-letter label is doing the work of a full experimental
specification, and the unreported remainder of the stack can move a result by
more than the algorithmic effect anyone is trying to claim. We have argued that
this stack-conditioning is the default, demonstrated it with a matched-config
comparison that flips by $17\times$, and proposed two coupled remedies: a
eight-item minimum-reportable-stack (\minreport{}), where every item is a
documented flip lever, and a controlled single-stack audit of DAPO, GSPO,
Dr.GRPO, and MAD-GRPO that reports which claimed gains survive. Neither remedy
requires new science---only a few fields of telemetry and the discipline to log
them. The cost is a JSON block; the payoff is a literature whose comparisons one
can actually trust.

% =============================================================================
% Bibliography. MANUAL STUB ONLY. We do NOT invent authors/years. Every entry
% is a TODO key to be replaced by a real citation before submission. When
% promoting to a venue template, delete this block, use \usepackage{natbib},
% and point \bibliography at a real .bib.
%
% SUBMISSION GATE: a reviewer of this manuscript will see a red
% [cite:KEY] marker for every in-text citation and a red [TODO] marker
% for every entry below. The paper will be treated as a draft, not as
% a submission, until the bibliography is real. The 12 cited works are
% all findable; if the position paper cannot stand on the citations it
% has, the claims need to be cut, not the citations stubbed.
% =============================================================================
\section*{References}
\small
\noindent\textcolor{red}{\textbf{[TODO: replace this stub with a real .bib.
Every \texttt{[cite:KEY]} marker above corresponds to one entry below. Do not
ship with placeholder keys. Each entry below carries a [\textsc{required}]
tag indicating its status; all are required for submission.]}}
\begin{itemize}[leftmargin=1.4em, itemsep=1pt, label={}]
  \item \texttt{zvfaudit2026} --- \todo{the ZVF Program v1 audit paper (this group's own work: ``Reward Contrast, Not Algorithm Labels''); cite the arXiv/venue version once assigned.}
  \item \texttt{henderson2018deeprl} --- \todo{Henderson et al., ``Deep Reinforcement Learning that Matters'' (reproducibility crisis in deep RL).}
  \item \texttt{islam2017reproducibility} --- \todo{Islam et al., reproducibility of benchmarked deep RL (verify exact reference).}
  \item \texttt{dapo2025} --- \todo{DAPO paper (clip-higher + dynamic sampling GRPO variant); fill authors/venue/year.}
  \item \texttt{gspo2025} --- \todo{GSPO paper (sequence-level importance sampling); fill authors/venue/year.}
  \item \texttt{drgrpo2025} --- \todo{Dr.GRPO paper (length/normalization fix); fill authors/venue/year.}
  \item \texttt{madgrpo2025} --- \todo{MAD-GRPO paper (multi-agent/diversity GRPO); fill authors/venue/year.}
  \item \texttt{aero2024} --- \todo{AERO (adaptive rollout sizing); from variance-mitigation bib fragment.}
  \item \texttt{cppo2024} --- \todo{CPPO (clip-pruned PPO); from variance-mitigation bib fragment.}
  \item \texttt{ngrpo2025} --- \todo{NGRPO (normalized GRPO); from variance-mitigation bib fragment.}
  \item \texttt{scafgrpo2025} --- \todo{Scaf-GRPO (scaffolded exploration); from variance-mitigation bib fragment.}
  \item \texttt{dar2024} --- \todo{DAR (dual-alignment / dual-KL regularization); from extended-RW bib fragment.}
\end{itemize}

\end{document}
